\section{总结与展望}

\subsection{论文工作总结}

% 请在此处总结论文工作

本文针对芳樟醇传统生产方法存在的问题，开展了多酶级联催化绿色合成芳樟醇的研究。本文的主要工作总结如下：

\textbf{（1）成功筛选并表达了芳樟醇合成的关键酶}

% 总结第一个贡献

从薄荷和薰衣草中筛选了高活性的GPPS和LIS基因，通过基因工程手段在大肠杆菌中实现了高效异源表达。

\textbf{（2）构建了高效的多酶级联催化体系}

% 总结第二个贡献

构建了GPPS和LIS的双酶级联反应体系，通过系统优化反应条件，芳樟醇的摩尔转化率达到85.3\%。

\textbf{（3）开发了固定化酶技术提高酶的稳定性}

% 总结第三个贡献

采用磁性纳米粒子固定化技术制备了固定化酶，显著提高了酶的稳定性和重复使用性能。

\textbf{（4）建立了芳樟醇绿色合成的技术路线}

% 总结第四个贡献

本研究建立了一条从简单前体IPP和DMAPP出发，通过多酶级联催化合成芳樟醇的绿色技术路线。

\subsection{研究局限性和未来展望}

% 请在此处讨论局限性和未来方向

尽管本研究取得了一定的成果，但仍存在一些局限性。

\textbf{（1）当前研究的局限性}

\textit{局限性一：底物成本较高}。

% 描述第一个局限性

本研究使用的底物IPP和DMAPP价格较高，限制了该方法的经济可行性。

\textit{局限性二：酶活性有待进一步提高}。

% 描述第二个局限性

目前GPPS和LIS的催化效率仍有提升空间。

\textbf{（2）未来研究方向}

基于上述局限性，未来研究可以从以下几个方向展开：

\textit{方向一：构建全细胞催化体系}。

% 描述第一个未来方向

将GPPS和LIS基因整合到同一宿主菌中，构建产芳樟醇的工程菌株。

\textit{方向二：酶的定向进化与改造}。

% 描述第二个未来方向

利用定向进化、理性设计等蛋白质工程技术，对GPPS和LIS进行改造。

\textit{方向三：反应体系的放大与工艺优化}。

% 描述第三个未来方向

开展反应体系的放大研究，建立适合工业化生产的生产工艺。
